\documentclass[aspectratio=169]{beamer}

\usetheme{Madrid}
\useinnertheme{circles}
\usecolortheme{default}

\usepackage[utf8]{inputenc}
\usepackage[T5]{fontenc}
\usepackage[vietnamese]{babel}
\usepackage{amsmath}
\usepackage{amssymb}
\usepackage{booktabs}
\usepackage{tikz}
\usetikzlibrary{arrows.meta, positioning}

\hypersetup{
  pdftitle={He thong kiem soat ra vao thong minh - AIoT},
  pdfauthor={Nhom: Quang, Nhan, Van, Phu},
  pdfsubject={Smart Access Control with Camera, RFID and Web Alerts},
}

\title[Kiểm soát ra vào AIoT]{Hệ thống kiểm soát ra vào thông minh}
\subtitle{Camera $\cdot$ RFID/NFC $\cdot$ Cảnh báo web}
\author[Nhóm thực hiện]{Lưu Huy Minh Quang -- 23127016 \and Hoàng Nhân -- 23127094 \and Trần Cao Vân -- 23127141 \and Nguyễn Trần Thiên Phú -- 23127246}
\institute[HCMUS]{Khoa Công nghệ Thông tin\\Đại học Khoa học Tự nhiên, ĐHQG TP.HCM\\[0.3em]\textit{Môn: Nhập môn lập trình điều khiển thiết bị thông minh}}
\date{Tháng 6, 2026}

\AtBeginSection[]{
  \begin{frame}[plain]
    \vfill
    \centering
      {\Large Phần \insertsectionnumber}\par\vspace{0.6em}
      {\LARGE\bfseries \insertsectionhead}\par
    \vfill
  \end{frame}
}

\begin{document}

%------- TITLE -------
\begin{frame}
  \begin{center}
    \vspace{0.4em}
    {\usebeamerfont{title}\Large\bfseries\inserttitle\par}
    \vspace{0.3em}
    {\usebeamerfont{subtitle}\insertsubtitle\par}

    \vspace{1.2em}
    {\large\bfseries Nhóm thực hiện}\\[0.6em]
    \begin{tabular}{ll}
      Lưu Huy Minh Quang & 23127016 \\
      Hoàng Nhân & 23127094 \\
      Trần Cao Vân & 23127141 \\
      Nguyễn Trần Thiên Phú & 23127246 \\
    \end{tabular}\\[0.8em]
    \textit{Khoa CNTT -- ĐH Khoa học Tự nhiên, ĐHQG TP.HCM}\\
    \textit{Môn: Nhập môn lập trình điều khiển thiết bị thông minh}

    \vspace{0.8em}
    {\large \insertdate}
  \end{center}
\end{frame}

%------- MỤC LỤC -------
\begin{frame}{Nội dung trình bày}
  \vspace{0.4em}
  \begin{enumerate}
    \item Giới thiệu project
    \item Mục đích của project
    \item Mục tiêu của project
    \item Phương pháp đề xuất
    \item Kiến trúc hệ thống
  \end{enumerate}
\end{frame}

%============================
\section{Giới thiệu project}
%============================

\begin{frame}{Nhóm đang làm gì?}
  \textbf{Sản phẩm:} Nhóm xây dựng một \textbf{hệ thống kiểm soát ra vào thông minh} giúp quản lý cửa ra vào an toàn và thuận tiện hơn. Hệ thống kết hợp \textbf{camera}, \textbf{RFID/NFC}, \textbf{servo} và \textbf{web dashboard} để nhận diện người, xác thực mở cửa và gửi cảnh báo khi có xâm nhập.

  \vspace{0.8em}
  \begin{itemize}
    \item Camera hỗ trợ phát hiện và nhận diện người đứng trước cửa.
    \item RFID/NFC giúp xác thực nhanh cho người dùng hợp lệ.
    \item Servo thực hiện đóng/mở cửa tự động sau khi xác thực.
    \item Hệ thống web nhận log ra vào và cảnh báo sự cố.
  \end{itemize}
\end{frame}

\begin{frame}{Dùng cho việc gì? -- Bối cảnh}
  \textbf{Bối cảnh sử dụng:} hệ thống phù hợp với các khu vực cần kiểm soát ra vào như phòng lab, văn phòng nhỏ, phòng server hoặc khu vực nội bộ.

  \vspace{0.8em}
  \begin{itemize}
    \item Quy mô sử dụng khoảng \textbf{20--30 người}.
    \item Cần quản lý người ra vào rõ ràng hơn.
    \item Cần hỗ trợ theo dõi và cảnh báo khi có bất thường.
  \end{itemize}
\end{frame}

\begin{frame}{Dùng cho việc gì? -- Hạn chế hiện tại}
  \textbf{Hạn chế của cách quản lý thủ công}

  \vspace{0.8em}
  \begin{itemize}
    \item Dễ xảy ra sai sót hoặc bỏ sót khi kiểm soát ra vào.
    \item Khó theo dõi đầy đủ lịch sử ra vào của từng người.
    \item Khó phản ứng kịp thời khi có người vào không đúng quy trình.
  \end{itemize}
\end{frame}

\begin{frame}{Dùng cho việc gì? -- Giải pháp đề xuất}
  \textbf{Giải pháp của nhóm}

  \vspace{0.8em}
  \begin{itemize}
    \item Xác thực người dùng nhanh hơn tại cửa.
    \item Ghi nhận lịch sử ra vào tự động trên hệ thống web.
    \item Cảnh báo sớm khi phát hiện vi phạm hoặc hành vi bất thường.
  \end{itemize}
\end{frame}

%============================
\section{Mục đích của project}
%============================

\begin{frame}{Mục đích}
  Sản phẩm được xây dựng nhằm tạo ra một hệ thống kiểm soát ra vào thông minh, giúp nhận diện hoặc xác thực người dùng hợp lệ để mở cửa tự động, đồng thời phát hiện xâm nhập trái phép và gửi cảnh báo nhanh cho người quản lý.
\end{frame}

%============================
\section{Mục tiêu của project}
%============================

\begin{frame}{Mục tiêu 1 -- Phát hiện và nhận diện người}
  \textbf{Human Detection \& Recognition}

  \vspace{0.8em}
  Hệ thống dùng camera để phát hiện và nhận diện người đứng trước cửa trong phạm vi \textbf{$0.5$ m}, với độ chính xác \textbf{trên $95\%$} và thời gian phản hồi \textbf{dưới $1.5$ giây}.
\end{frame}

\begin{frame}{Mục tiêu 2 -- Phát hiện trèo/nhảy qua cổng}
  \textbf{Gate Intrusion Detection}

  \vspace{0.8em}
  Hệ thống theo dõi khu vực cổng để phát hiện người trèo hoặc nhảy qua trong vòng \textbf{$2$ giây}, sau đó báo động tại chỗ và gửi cảnh báo lên hệ thống web để người quản lý xử lý kịp thời.
\end{frame}

\begin{frame}{Mục tiêu 3 -- Xác thực RFID/NFC và mở cửa tự động}
  \textbf{RFID/NFC Authentication \& Auto Unlock}

  \vspace{0.8em}
  Hệ thống dùng RFID/NFC để xác thực người dùng hợp lệ và mở cửa tự động trong vòng \textbf{$3$ giây}, giúp việc ra vào nhanh và thuận tiện hơn.
\end{frame}

%============================
\section{Phương pháp đề xuất}
%============================

\begin{frame}{Phương pháp thực hiện}
  Hệ thống được triển khai theo quy trình đơn giản: nhận tín hiệu từ camera, RFID/NFC và cảm biến; xử lý tại bộ điều khiển; sau đó mở cửa hoặc kích hoạt cảnh báo và cập nhật trạng thái lên web.

  \vspace{0.8em}
  \begin{center}
  \begin{tikzpicture}[
    pipestep/.style={rectangle, rounded corners=4pt, draw=blue!70, fill=blue!8,
      text width=2.8cm, align=center, font=\small, minimum height=1.0cm},
    arr/.style={-{Stealth[length=2.5mm]}, thick, blue!70}
  ]
    \node[pipestep, fill=teal!8, draw=teal!70!black] (input) {Nhận tín hiệu\\{\scriptsize Camera / RFID / cảm biến}};
    \node[pipestep, right=0.7cm of input] (process) {Xử lý tại thiết bị\\{\scriptsize kiểm tra dữ liệu}};
    \node[pipestep, right=0.7cm of process] (decision) {Ra quyết định\\{\scriptsize hợp lệ hay vi phạm}};
    \node[pipestep, right=0.7cm of decision, fill=orange!10, draw=orange!80!black] (output) {Hành động\\{\scriptsize mở cửa / báo động / gửi web}};

    \draw[arr] (input) -- (process);
    \draw[arr] (process) -- (decision);
    \draw[arr] (decision) -- (output);
  \end{tikzpicture}
  \end{center}
\end{frame}

\begin{frame}{Luồng 1 -- Nhận diện bằng camera}
  Camera được đặt trước cửa để ghi nhận hình ảnh người đến gần. Dữ liệu ảnh được xử lý tại thiết bị nhằm phát hiện có người đứng trước cửa và hỗ trợ nhận diện người dùng đã đăng ký trong hệ thống.

  \vspace{0.8em}
  \begin{itemize}
    \item Phát hiện người trong phạm vi hoạt động yêu cầu.
    \item Đối chiếu dữ liệu nhận diện với thông tin đã lưu.
    \item Gửi kết quả sang bộ điều khiển để quyết định mở cửa hoặc từ chối.
  \end{itemize}
\end{frame}

\begin{frame}{Luồng 2 -- Xác thực RFID/NFC}
  Người dùng hợp lệ có thể sử dụng thẻ RFID/NFC để xác thực nhanh khi ra vào. Sau khi quét thẻ, bộ điều khiển kiểm tra thông tin và thực hiện mở cửa tự động nếu dữ liệu hợp lệ.

  \vspace{0.8em}
  \begin{itemize}
    \item Đọc mã thẻ từ đầu đọc RFID/NFC.
    \item Kiểm tra mã thẻ với danh sách người dùng đã đăng ký.
    \item Mở cửa và ghi nhận log khi xác thực thành công.
  \end{itemize}
\end{frame}

\begin{frame}{Luồng 3 -- Phát hiện trèo/nhảy qua cổng}
  Hệ thống theo dõi khu vực cổng bằng camera hoặc cảm biến để nhận biết các hành vi xâm nhập bất thường như trèo hoặc nhảy qua cổng mà không thực hiện xác thực hợp lệ.

  \vspace{0.8em}
  \begin{itemize}
    \item Theo dõi khu vực cổng liên tục trong thời gian hoạt động.
    \item Phát hiện nhanh các trường hợp vượt cổng trái phép.
    \item Kích hoạt báo động và gửi cảnh báo lên web cho người quản lý.
  \end{itemize}
\end{frame}

\begin{frame}{Web quản lý và cảnh báo}
  Hệ thống web đóng vai trò theo dõi và quản lý tập trung. Người quản lý có thể xem thông tin người dùng, lịch sử ra vào, trạng thái cửa và các cảnh báo phát sinh từ thiết bị.

  \vspace{0.8em}
  \begin{itemize}
    \item Quản lý danh sách người dùng và thẻ RFID/NFC.
    \item Hiển thị log ra vào theo thời gian thực.
    \item Nhận cảnh báo khi phát hiện trèo/nhảy qua cổng.
    \item Hỗ trợ theo dõi trạng thái hoạt động của hệ thống.
  \end{itemize}
\end{frame}

\begin{frame}{Luồng hoạt động tổng quát}
  \begin{enumerate}
    \item Khi có người đến cửa, hệ thống nhận tín hiệu từ camera hoặc RFID/NFC.
    \item Bộ điều khiển xử lý dữ liệu để xác định người dùng hợp lệ hay trường hợp bất thường.
    \item Nếu hợp lệ, servo mở cửa và hệ thống lưu lại log ra vào.
    \item Nếu phát hiện trèo/nhảy qua cổng, hệ thống bật báo động và gửi cảnh báo lên web.
  \end{enumerate}
\end{frame}

\begin{frame}{Ưu điểm}
  \begin{itemize}
    \item \textbf{Ra vào thuận tiện:} hỗ trợ nhận diện người và xác thực bằng RFID/NFC, giúp mở cửa nhanh hơn.
    \vspace{0.4em}
    \item \textbf{Quản lý dễ dàng:} toàn bộ lịch sử ra vào và trạng thái hệ thống được theo dõi trên web.
    \vspace{0.4em}
    \item \textbf{An toàn hơn:} hệ thống có thể phát hiện vượt cổng và cảnh báo ngay tại chỗ lẫn trên web.
  \end{itemize}
\end{frame}

\begin{frame}{Các khó khăn có thể gặp phải}
  \begin{itemize}
    \item \textbf{Môi trường thực tế:} ánh sáng, góc đứng và khoảng cách có thể ảnh hưởng đến chất lượng nhận diện bằng camera.
    \vspace{0.3em}
    \item \textbf{Độ ổn định phần cứng:} đầu đọc RFID/NFC, cảm biến và servo cần hoạt động đồng bộ để tránh lỗi khi vận hành.
    \vspace{0.3em}
    \item \textbf{Tích hợp hệ thống:} cần đồng bộ tốt giữa thiết bị, web và cảnh báo để bảo đảm phản hồi đúng lúc.
  \end{itemize}
\end{frame}

%============================
\section{Kiến trúc hệ thống}
%============================

\begin{frame}{Kiến trúc tổng thể}
  \begin{center}
  \begin{tikzpicture}[
    layerbox/.style={rectangle, rounded corners=4pt, text width=3.2cm,
      align=center, minimum height=1.5cm, font=\small\bfseries, inner sep=5pt},
    inputbox/.style={layerbox, draw=blue!70, fill=blue!8},
    ctrlbox/.style={layerbox, draw=violet!70!black, fill=violet!10},
    actbox/.style={layerbox, draw=green!55!black, fill=green!9},
    webbox/.style={layerbox, draw=orange!80!black, fill=orange!10},
    arr/.style={-{Stealth[length=2mm]}, thick, blue!70}
  ]
    \node[inputbox] (input) {Khối đầu vào\\{\scriptsize Camera + RFID/NFC + cảm biến cổng}};
    \node[ctrlbox, right=1.0cm of input] (ctrl) {Khối xử lý trung tâm\\{\scriptsize ESP32 / ESP32-CAM}};
    \node[actbox, right=1.0cm of ctrl] (act) {Khối chấp hành\\{\scriptsize Servo + LED/Buzzer}};
    \node[webbox, above=1.3cm of ctrl] (web) {Web dashboard\\{\scriptsize log + trạng thái + cảnh báo}};

    \draw[arr] (input) -- (ctrl);
    \draw[arr] (ctrl) -- (act);
    \draw[arr] (ctrl) -- (web);
  \end{tikzpicture}
  \end{center}

  \vspace{0.4em}
  \textbf{Cách hoạt động:} khối đầu vào nhận tín hiệu, khối xử lý trung tâm ra quyết định, khối chấp hành thực hiện mở cửa hoặc báo động, còn web dùng để theo dõi và quản lý.
\end{frame}

\begin{frame}{Kiến trúc -- Chi tiết thành phần}
  \begin{columns}[T]
    \column{0.33\textwidth}
      \textbf{Khối đầu vào}
      \begin{itemize}
        \item Camera trước cửa
        \item Đầu đọc RFID/NFC
        \item Cảm biến phát hiện vượt cổng
        \item Tín hiệu người dùng và môi trường
      \end{itemize}

    \column{0.33\textwidth}
      \textbf{Khối xử lý}
      \begin{itemize}
        \item ESP32 / ESP32-CAM
        \item Xử lý nhận diện và xác thực
        \item So khớp dữ liệu người dùng
        \item Ra quyết định mở cửa hoặc cảnh báo
      \end{itemize}

    \column{0.33\textwidth}
      \textbf{Khối đầu ra \& giám sát}
      \begin{itemize}
        \item Servo đóng/mở cửa
        \item LED / Buzzer báo trạng thái
        \item Web dashboard
        \item Log ra vào và cảnh báo
      \end{itemize}
  \end{columns}
\end{frame}


\begin{frame}{Vòng lặp điều khiển thông minh}
  \begin{center}
  \begin{tikzpicture}[
    box/.style={rectangle, rounded corners=3pt, draw=blue!70, fill=blue!7,
      text width=2.5cm, align=center, font=\small, minimum height=0.9cm},
    arr/.style={-{Stealth[length=2mm]}, thick, blue!70}
  ]
    \node[box, fill=teal!8, draw=teal!70!black] (sense) {Cảm nhận\\{\scriptsize Camera / RFID / cảm biến}};
    \node[box, right=0.6cm of sense] (detect) {Xử lý\\{\scriptsize kiểm tra dữ liệu}};
    \node[box, right=0.6cm of detect] (process) {Quyết định\\{\scriptsize hợp lệ hay vi phạm}};
    \node[box, right=0.6cm of process, fill=orange!10, draw=orange!80!black] (decide) {Hành động\\{\scriptsize mở cửa / báo động}};
    \node[box, below=0.8cm of decide, fill=green!9, draw=green!55!black] (act) {Giám sát\\{\scriptsize ghi log lên web}};

    \draw[arr] (sense) -- (detect);
    \draw[arr] (detect) -- (process);
    \draw[arr] (process) -- (decide);
    \draw[arr] (decide) -- (act);
    \draw[-{Stealth[length=2mm]}, thick, red!60] (act.west) -| node[below, font=\scriptsize, red!60, near start]{Phản hồi trạng thái} (sense.south);
  \end{tikzpicture}
  \end{center}

  \vspace{0.6em}
  \textbf{Ý tưởng chính:} hệ thống liên tục nhận tín hiệu, xử lý, ra quyết định, thực hiện hành động và cập nhật trạng thái để người quản lý theo dõi.
\end{frame}

\begin{frame}{Danh mục linh kiện dự kiến}
  \scriptsize
  \renewcommand{\arraystretch}{1.3}
  \begin{tabular}{p{0.4cm}p{2.5cm}p{6.2cm}p{1.0cm}p{1.5cm}}
    \toprule
    \textbf{STT} & \textbf{Linh kiện} & \textbf{Chi tiết kỹ thuật} & \textbf{SL} & \textbf{Giá (VNĐ)} \\
    \midrule
    1 & ESP32-38 pin & Bộ điều khiển trung tâm, Wi-Fi 2.4GHz, dùng để nhận tín hiệu và điều khiển thiết bị & 1 & $\sim$260.000 \\
    2 & ESP32-S3-CAM & Camera OV2640, dùng để thu hình và hỗ trợ xử lý nhận diện & 1 & $\sim$280.000 \\
    3 & Module RFID/NFC & RC522, dùng để đọc thẻ xác thực người dùng & 1 & $\sim$45.000 \\
    4 & Cảm biến khoảng cách / hồng ngoại & Dùng để phát hiện người đến gần hoặc vượt cổng & 1 & $\sim$35.000 \\
    5 & Động cơ Servo & SG90, dùng để mô phỏng cơ chế đóng/mở cửa & 2 & $\sim$60.000 \\
    6 & Breadboard & 170 lỗ cắm, phục vụ lắp ráp và thử nghiệm mạch & 1 & $\sim$60.000 \\
    7 & Phụ kiện và nguồn & Dây jump, LED, Buzzer, điện trở, adapter cấp nguồn & 1 bộ & $\sim$120.000 \\
    \midrule
    \multicolumn{4}{r}{\textbf{Tổng cộng (ước tính):}} & \textbf{$\sim$860.000} \\
    \bottomrule
  \end{tabular}
\end{frame}

%------- CLOSING -------
\begin{frame}[plain]
  \begin{center}
    \vspace{2.5cm}
    {\Huge \textbf{Cảm ơn!}}\\[0.8em]
    {\Large Câu hỏi \& Thảo luận}
  \end{center}
\end{frame}

\end{document}
